\subsection{Partielle Ableitungen und Richtungsableitungen}

Sei $ D \subset \mathbb{R}^n $ eine offene Menge und $ f: D \to \mathbb{R} $ eine Funktion. Wir betrachten Ableitungen von $ f $ in einzelnen Richtungen und definieren damit auch die partiellen Ableitungen, indem wir zunächst die Richtungsableitungen einführen.

\begin{definition}{}
1) Ist $ f $ in einem Punkt $ a \in D $ (im Sinne einer eindimensionalen Differenzierbarkeit in der Richtung) differenzierbar, so heißt die Richtungsableitung von $ f $ in $ a $ in Richtung $ v \in \mathbb{R}^n $


$$
D_v f(a) = \lim_{t\to 0} \frac{f(a+tv)-f(a)}{t},
$$


vorausgesetzt, dieser Grenzwert existiert. (Hinweis: Auch für den Fall $ v=0 $ ist diese Definition formell sinnvoll – es ergibt sich dabei trivialerweise $ D_{0}f(a)=0 $.)


2) $ f $ heißt an der Stelle $ a $ \emph{partiell differenzierbar}, wenn für alle Richtungen der Standardbasis, also insbesondere für die Richtungsableitungen in den Richtungen $ e_1,\dots,e_n $, der jeweilige Grenzwert existiert. Wir setzen dann


$$
\frac{\partial f}{\partial x_i}(a) = D_{e_i} f(a), \quad i=1,\dots,n.
$$


3) $ f $ heißt auf $ D $ (stetig) partiell differenzierbar, wenn sie in jedem Punkt $ a \in D $ partiell differenzierbar ist.
\end{definition}

\begin{bemerkungbox}
    Die partiellen Ableitungen $\frac{\partial f}{\partial x_i}(a)$ messen die Änderungsrate von $ f $ in den Richtungen der Koordinatenachsen $ e_i $. Sie liefern also die "Ableitung“ der Funktion entlang der einzelnen Koordinaten.
\end{bemerkungbox}

\begin{beispielbox}
\begin{enumerate}
  \item \textbf{Funktion in einer Variablen:} Sei $ f(x)=x^2 $ für $ x\in \mathbb{R} $. Dann existiert die (gewöhnliche) Ableitung
  

$$
  \frac{d}{dx} f(x)= 2x.
  $$


  
  \item \textbf{Funktion in zwei Variablen:} Betrachte $ f:\mathbb{R}^2 \to \mathbb{R} $ definiert durch
  

$$
  f(x,y)= x^2+y^2.
  $$


  Dann sind die partiellen Ableitungen
  

$$
  \frac{\partial f}{\partial x}(x,y)=2x,\quad \frac{\partial f}{\partial y}(x,y)=2y.
  $$


  Insbesondere gilt beispielsweise für $ y=0 $ natürlich
  

$$
  \frac{\partial f}{\partial y}(x,0)=0.
  $$


  
  \item \textbf{Hinweis zu Randfällen:} In manchen Fällen können sich Richtungsableitungen in anderen Richtungen ($ v\notin\{e_1,\dots,e_n\} $) anders verhalten – oftmals geht es im Detail um die Frage, ob alle Richtungsableitungen existieren und ob diese in konzisen Situationen übereinstimmen.
\end{enumerate}
\end{beispielbox}

\begin{bemerkungbox}
Kurzer Einschub: Ist $ D_v f(a) = \lim_{t\to 0} \frac{f(a+tv)-f(a)}{t} $ für eine beliebige Richtung $ v $ definiert, so gibt diese Richtungsableitung die Steigung der Tangente an den Funktionsgraphen in Richtung $ v $ an. Im speziellen Fall $ n=1 $ entspricht dies der gewöhnlichen Ableitung. Allgemein wird so die geometrische Interpretation der partiellen Ableitungen verallgemeinert.
\end{bemerkungbox}


\subsection*{Gradient und Divergenz}

\begin{definition}{}
(a) Es sei $ f : D \to \mathbb{R} $ ($D \subset \mathbb{R}^n$) partiell differenzierbar an einem Punkt $ a \in D $. Dann heißt der Vektor


$$
    \nabla f(a) \equiv \operatorname{grad} f(a) := \Bigl( \frac{\partial f}{\partial x_1}(a), \frac{\partial f}{\partial x_2}(a), \dots, \frac{\partial f}{\partial x_n}(a) \Bigr) \in \mathbb{R}^n \qquad [\nabla \ = \ 'Nabla']
$$

der \emph{Gradient} von $ f $ in $ a $.
\newline

(b) Ein \emph{Vektorfeld} auf einer offenen Menge $ D \subset \mathbb{R}^n $ ist eine Abbildung


$$
v: D \to \mathbb{R}^n,\quad x \mapsto v(x).
$$



(c) Das Vektorfeld $ v $ heißt \emph{(stetig) partiell differenzierbar}, falls alle Komponentenfunktionen $ v_i: D \to \mathbb{R} $ (für $ i=1,\dots,n $) (stetig) partiell differenzierbar sind. In diesem Fall definieren wir die \emph{Divergenz} von $ v $ durch


$$
\operatorname{div} v : D \rightarrow \mathbb{R}, \quad
\operatorname{div} v(x) := \sum_{i=1}^n \frac{\partial v_i}{\partial x_i}(x).
$$


\end{definition}

\begin{bemerkungbox}
Üblicherweise schreibt man den Gradienten als $\nabla f$ und notiert dabei


$$
\langle \nabla f(a), e_i \rangle = \frac{\partial f}{\partial x_i}(a),
$$


wobei $(e_1, \dots, e_n)$ die Standardbasis von $\mathbb{R}^n$ ist.
\end{bemerkungbox}

\begin{beispielbox}
\begin{enumerate}
  \item \textbf{Notation als Differenzialoperator:} Man schreibt
  

$$
  \operatorname{grad} f = \nabla f.
  $$


  
  \item \textbf{Radiale Koordinate:} Sei $ r : \mathbb{R}^n \to \mathbb{R} $ definiert durch
  

$$
  r(x) = \|x\|,
  $$


  sodass $ r $ partiell differenzierbar ist (für $ x\neq 0 $). Dann gilt
  

$$
  \frac{\partial r}{\partial x_i}(x)= \frac{x_i}{\|x\|} \quad\text{und}\quad \operatorname{grad} r(x)= \frac{x}{\|x\|}.
  $$


  
  \item \textbf{Divergenz eines Vektorfeldes:} Sei $ v(x) = r(x)\, x $, wobei $ r(x) $ eine skalare Funktion ist. Dann folgt mit der Produktregel
  

$$
  \operatorname{div}\bigl( v(x) \bigr)= \langle \nabla r(x), x \rangle + r(x) \cdot \operatorname{div}(x),
  $$


  wobei $ \operatorname{div}(x)=n $ für $ x\in \mathbb{R}^n $ verwendet wird.
\end{enumerate}
\end{beispielbox}
